\section{Introducción}\label{sec:intro}

Un manipulador industrial heredado es un robot cuya mecánica sigue siendo útil, pero cuya controladora ha quedado obsoleta, es cerrada o carece de las interfaces de red que esperan los sistemas actuales. Sustituirlo no siempre es viable para una pyme o un centro de formación, y la modernización (en inglés \textit{retrofitting}) se ha consolidado como alternativa para incorporarlo a la Industria 4.0. En la revisión sistemática de \citet{alqoud2022industry}, con 32 trabajos publicados entre 2015 y 2021, los brazos robóticos son el segundo tipo de máquina más estudiado (16~\%), muy por detrás de las máquinas de control numérico (50~\%). En docencia el problema se repite con robots didácticos e industriales de los años 80 y 90 que conservan la mecánica intacta y tienen la electrónica averiada o cerrada \citep{conchasanchez2020recycling,roldan2018enhancing}.

En los manipuladores, \citet{bilancia2023overview} distinguen entre conservar la controladora original y gobernarla desde fuera, como coordinador externo (\textit{external coordinator}), y sustituirla total o parcialmente por un controlador abierto (\textit{open controller}). En los robots antiguos predomina la sustitución.\citet{alvares2018retrofitting} sustituyen la controladora de un ASEA IRB6-S2 por un control numérico abierto sobre LinuxCNC, \citet{jokic2020two} gobiernan un PUMA 560 desde un PC o una FPGA y \citet{conchasanchez2020recycling} actualizan un robot didáctico ED7220C con un Arduino y Simulink. En el propio RV-M1, \citet{roldan2018enhancing} sustituyen la controladora por una dSPACE DS1103 con servoamplificadores propios y obtienen un error en régimen permanente de 0,3~mm en una trayectoria.

Como contrapartida, abrir el lazo obliga a renunciar a la controladora original, a su planificación de trayectorias y a sus seguridades. Cuando la controladora se conserva, suele tener interfaz de red aunque sea heredada. Por ejemplo, \citet{holmes2025retrofitting} integran dos robots KUKA con controladoras KRC2 heredadas mediante una Raspberry Pi que se comunica con ellas por OPC XML-DA y un servidor OPC UA por dispositivo, y miden un 8,3 y un 15,4~\% de fallos en las operaciones de coger y dejar, pero no la latencia.

En el RV-M1, la vía no invasiva se exploró entre 2005 y 2008 a través del puerto RS-232 de la controladora. \citet{swider2007system} desarrollan un sistema de simulación y programación remota, \citet{foit2008web} lo lleva a una interfaz web y \citet{mihali2005effective} presentan un simulador OpenGL que se mueve a la par que el robot, en el que el usuario elige al conectar si gobierna el robot o el simulador. Los tres tienen una motivación docente, usan tecnologías hoy obsoletas y no aportan medidas. Más recientemente, \citet{patrao2025understanding} estudian integrar un Mitsubishi RV-2AJ, con su controladora y su RS-232, en el gemelo digital de un sistema de fabricación flexible heredado, pero dejan sin implementar la realimentación desde la controladora. Fuera de Mitsubishi, la ScorBot Toolbox de \citet{kutzer2023scorbot} gobierna un SCORBOT-ER 4u con su controladora original y ofrece una simulación con la misma sintaxis que las funciones del robot. Es código abierto, pero funciona en local y no mide la discrepancia entre la simulación y el robot.

Por otro lado, MQTT se usa ya como capa de mensajería entre los manipuladores y sus gemelos digitales. En un banco de pruebas en red local, las diferencias de tiempo de ida y vuelta entre MQTT y OPC UA cliente/servidor fueron inferiores al milisegundo \citep{profanter2019opcua}, mientras que en una celda con un UR5e MQTT fue varios milisegundos más rápido que OPC UA PubSub, con latencias medianas de lazo cerrado de 2,5 y 10,1~ms \citep{ariascachero2023operational}. MQTT se ha empleado también en teleoperación con distintos niveles de calidad de servicio (QoS) y cifrado TLS \citep{luu2025enhancing}. Sin embargo, todas estas medidas corresponden a controladoras con interfaz de red, la latencia rara vez se desglosa por tramos \citep{bray2024latency} y no se han encontrado datos del retardo que añaden un enlace serie y el intérprete de órdenes de una controladora de los años 80.

En los gemelos digitales de manipuladores el acoplamiento con el robot varía mucho. \citet{kaarlela2023digital} dejan elegir al alumno entre modelo digital, sombra digital y gemelo digital, con MQTT como protocolo principal y un puente OPC UA/MQTT hacia la controladora. \FALTA{cita de definición de modelo, sombra y gemelo digital (p.~ej., Kritzinger et al., 2018), pendiente de añadir a refs.bib y de verificar.} Los gemelos bidireccionales se apoyan en el estado que ofrece una controladora con interfaz de red, de forma continua o por sondeo frecuente, como el de \citet{chang2025construction} sobre un FANUC a 125~Hz, el de \citet{liang2020bidirectional} con un puente MQTT a 250~Hz o el de \citet{singh2024unity} sobre ROS.

Que un mismo programa gobierne el robot real o su simulación tampoco es nuevo. Es práctica habitual en ROS, cuya capa de abstracción de hardware permite intercambiar el robot real y el simulado \FALTA{cita de ros\_control (Chitta et al., 2017), pendiente de añadir y de verificar}, y aparece también en robots didácticos y de bajo coste \citep{mihali2005effective,kutzer2023scorbot,zhang2024construction}. Del mismo modo, la emulación de un dispositivo para que el software de control no lo distinga del real se ha formalizado como prototipo de gemelo digital (\textit{digital twin prototype}) \FALTA{cita de Barbie, Hasselbring y Hansen (2024), pendiente de añadir y de verificar}. Además, la fidelidad del gemelo se informa de forma heterogénea: en milímetros del punto central de la herramienta (\textit{tool centre point}, TCP) en pocos ensayos \citep{chang2025construction,zhang2024construction}, en radianes por articulación \citep{liang2020bidirectional} o como un porcentaje normalizado por 180\textdegree{} \citep{singh2024unity}. En ningún caso se separa la coherencia del gemelo con lo que informa la controladora de su fidelidad física frente a una medida externa.

En cuanto a la comunicación, varias salvaguardas útiles entre un gemelo, una pasarela y un robot están ya normalizadas. MQTT 3.1.1 define la última voluntad (\textit{last will and testament}, LWT), la entrega con QoS~1, que puede duplicar mensajes, y el indicador \textit{retain} \citep{oasis2014mqtt}. MQTT 5.0 añade una caducidad de los mensajes, expresada en segundos, que aplica el broker \citep{oasis2019mqtt}, y Sparkplug, normalizado como ISO/IEC 20237, gestiona sobre MQTT el nacimiento y la muerte de los nodos \citep{iso2023sparkplug}. \VERIFICAR{mecanismo de Sparkplug tomado de documentación secundaria; comprobar en el texto de ISO/IEC 20237.} Sin embargo, en los trabajos sobre manipuladores las salvaguardas se limitan al cifrado, la red privada virtual (VPN), los permisos y los umbrales de desincronización \citep{luu2025enhancing,bravoarrabal2025iort,liang2020bidirectional}, y la tolerancia a retrasos que declaran \citet{chang2025construction} no se describe ni se ensaya.

Por tanto, en la búsqueda realizada no se ha encontrado ningún trabajo que conserve una controladora industrial heredada cuya única interfaz sea un puerto serie, la conecte a un protocolo del Internet industrial de las cosas (\textit{Industrial Internet of Things}, IIoT) con un gemelo digital y caracterice el conjunto con medidas. La \cref{tab:comparativa} resume las diferencias con los trabajos más próximos. A diferencia de las controladoras que alimentan los gemelos revisados, la D/U-M1 del RV-M1 solo informa de su posición y de su estado de error cuando se le piden, con las órdenes \texttt{WH} y \texttt{ER}, a través de la línea RS-232C.

\begin{table}[t]
  \centering
  \caption{Comparación con trabajos representativos. Interfaz: canal entre la pasarela o el PC y la controladora. n.~i.: no indicado; n.~a.: no aplica.}
  \label{tab:comparativa}
  \footnotesize
  \setlength{\tabcolsep}{3pt}
  \begin{tabular}{@{}p{2.1cm}p{1.3cm}p{1.8cm}p{1.9cm}p{2.4cm}p{2.5cm}p{2.5cm}@{}}
    \toprule
    Trabajo (robot) & Controladora original & Interfaz & Protocolo & Gemelo digital & Seguridades & Medidas \\
    \midrule
    \citet{roldan2018enhancing} (RV-M1) & No (dSPACE) & E/S analógicas y digitales & n.~a. & No (simulación en Simulink) & n.~i. & Error en régimen de 0,3~mm (una trayectoria) \\
    \citet{holmes2025retrofitting} (2 KUKA) & Sí (KRC2) & OPC XML-DA (Ethernet) & OPC UA & n.~i. & n.~i. & Fallos de coger y dejar; tiempo perdido; sin latencia \\
    \citet{mihali2005effective} (RV-M1) & Sí & RS-232 & TCP/IP propio & Visualización sincronizada; robot o simulador al conectar & n.~i. & Ninguna \\
    \citet{kutzer2023scorbot} (SCORBOT-ER 4u) & Sí & USB (DLL del fabricante) & HTTP/JSON local & Simulación con la misma sintaxis, no sincronizada & n.~a. (local) & Métricas de código; valoración cualitativa \\
    \citet{patrao2025understanding} (RV-2AJ) & Sí & RS-232 & n.~i. & Propuesto & n.~i. & Ninguna \\
    \citet{kaarlela2023digital} (UR, Fanuc, ABB) & Sí & Red (OPC UA) & MQTT, WebSocket & Modelo, sombra y gemelo & n.~i. & Validación funcional \\
    \citet{liang2020bidirectional} (KUKA KR120) & Sí & Red (PLC) & ROS y puente MQTT & Bidireccional, 250~Hz & Umbral de pose & Error articular (rad); sin latencia \\
    \citet{chang2025construction} (FANUC) & Sí & Ethernet (librería propietaria) & TCP/IP & Bidireccional, sondeo a 125~Hz & Umbral de desincronización; parada de emergencia & Error XYZ de $\pm 1$~mm; sin latencia \\
    \citet{luu2025enhancing} (ABB IRB1200) & Sí & EGM/UDP & MQTT (QoS 0/1/2) & Visualización sincronizada (VR) & TLS y autenticación & Latencia por QoS; pérdidas \\
    \midrule
    Este trabajo (RV-M1) & Sí, sin modificar & RS-232C & MQTT (QoS 1, JSON) & Emulador, sombra y sincronizado; estado bajo petición (\texttt{WH}/\texttt{ER}) & LWT, caducidad, duplicados, destino doble, ACK en dos fases, ACL & Latencia por tramos; coherencia y fidelidad física; inyección de fallos \\
    \bottomrule
  \end{tabular}
\end{table}

El objetivo de este trabajo es presentar la modernización no invasiva de un Mitsubishi Movemaster EX RV-M1 cuya controladora D/U-M1 se conserva sin modificar, con su planificación de trayectorias y sus seguridades, y se gobierna por su RS-232C desde Python, desde un teach pendant de software y a través de un puente RS-232C/MQTT. En la taxonomía de \citet{bilancia2023overview} es, por tanto, un coordinador externo. El gobierno del RV-M1 por RS-232 desde software no es nuevo \citep{swider2007system,foit2008web,mihali2005effective}, y aquí es solo la base sobre la que se construyen las aportaciones.

La aportación principal es (i) un gemelo digital web con tres modos: emulador, en el que interpreta las órdenes del contrato MQTT como la D/U-M1 y responde con los mismos acuses y estados que el puente; sombra, en el que además sigue la pose que informa el robot real y calcula el error entre ambos; y sincronizado, en el que reproduce a distancia la pose del robot real. El gemelo aplica la emulación de dispositivos a una controladora que solo informa bajo petición, y su fidelidad se mide en milímetros y grados, separando la coherencia con la posición que devuelve la controladora y la fidelidad física frente a una medida externa. Las otras dos son (ii) la adaptación a una controladora serie sin interfaz de red de salvaguardas conocidas de MQTT y Sparkplug (última voluntad, caducidad comprobada también antes de escribir en la línea serie, descarte de duplicados, destino en el topic y en el mensaje y acuse en dos fases), junto con su verificación por inyección de fallos, y (iii) la caracterización de la cadena serie heredada, con la latencia desglosada por tramos y la robustez ante cortes del broker y del enlace serie.

El resto del artículo se organiza de la siguiente forma. En la \cref{sec:modernisation} se describe el proceso de modernización, en la \cref{sec:mqtt} se presenta el protocolo de comunicación MQTT, en la \cref{sec:experiments} se recogen los experimentos y resultados y, finalmente, en la \cref{sec:conclusions} se exponen las conclusiones.
